% !PRINTED 73
\weberpage{73}

Dann müssen die $\gamma$ der Bedingung genügen:

\weq{1}{C_{-1}^{\gamma_{-1}} C_0^{\gamma_0} C_1^{\gamma_1} \dots C_\mu^{\gamma_\mu} \equiv 1 \, (\mathrm{mod}\ m').}

Diese Bedingung fordert, dass alle Indices mit Ausnahme des der Primzahl $q$
entsprechenden, den wir mit $\gamma$ bezeichnen, Null sein müssen. In Bezug auf
$\gamma$ ist aber zu unterscheiden, ob $q$ noch in $m'$ aufgeht oder nicht, d.\,h.
ob $q$ ein mehrfacher oder nur ein einfacher Factor von $m$ ist. Ist $q$ nicht
mehr in $m'$ enthalten, so enthält die Congruenz (1) gar keine Beschränkung für
$\gamma$, und $\gamma$ kann jeden Werth nach dem Modul $c$, d.\,h. jeden Werth
$0, 1, \dots q - 2$ annehmen.

Ist aber $q$ noch in $m'$ enthalten, also $m$ durch $q^\varkappa$ theilbar, und $\varkappa > 1$,
so fordert die Bedingung (1):
\[
C^\gamma \equiv 1 \, (\mathrm{mod}\ q^{\varkappa-1}),
\]
also
\weq{2}{\gamma \equiv 0 \, \left(\mathrm{mod}\ \dfrac{c}{q}\right),}
und $\gamma$ kann also jeden der Werthe
\[
0, \dfrac{c}{q}, \dfrac{2c}{q}, \dots \dfrac{(q-1)c}{q}
\]
erhalten. Im ersten Falle ist der Grad der Gruppe $Q$ gleich $q - 1$, im
zweiten gleich $q$.

Die Indices $\delta$ einer Zahl aus der Gruppe $P$ erhält man nach \S.\,12, 7.
und \S.\,16, (7) aus der Bedingung, dass für alle zulässigen $\gamma$
\weq{3}{\varepsilon_{-1}^{\gamma_{-1}\delta_{-1}} \varepsilon_0^{\gamma_0\delta_0} \varepsilon_1^{\gamma_1\delta_1} \dots \varepsilon_\mu^{\gamma_\mu\delta_\mu} = 1}
sein soll. Nach dem, was eben über die Indices $\gamma$ bewiesen ist, fordert
aber (3) nur das eine, dass der der Primzahl $q$ entsprechende Index $\delta$
durch $q$ oder durch $q - 1$ theilbar sein soll, je nachdem $q$ mehrmals oder
nur einmal in $m$ aufgeht. Wir heben also den Satz hervor:

7. \gesperrt{Die zu $Q$ reciproke Gruppe $P$ ist dadurch charakterisirt, dass
der $q$ entsprechende Index $\delta$ aller seiner Zahlen durch $q$ oder durch
$q - 1$ theilbar ist, je nachdem $q$ mehrmals oder nur einmal in $m$ aufgeht.}

Und daraus:

8. \gesperrt{Ein nicht primärer Theiler $\mathfrak{A}$ von $\mathfrak{N}$ ist
dadurch charakterisirt, dass in der reciproken Gruppe $\mathfrak{B}$}
%CONT>>
